\documentclass{article}
\usepackage{amsmath,amssymb}

\title{Classic Malbolge Word And Arithmetic Model}
\author{Malbolge Project}
\date{2026-07-26}

\begin{document}
\maketitle

\section{Scope}

This document records normative mathematics for the classic 1998 Malbolge word
model. It does not define loader validation, instruction translation tables,
I/O, or extended Malbolge profiles.

\section{Word domain}

A classic machine word contains ten ternary digits. Let
\[
  W = 3^{10} = 59049,
  \qquad
  \mathcal{W} = \{0,1,\ldots,W-1\}.
\]
Every $x\in\mathcal{W}$ has a unique ten-trit expansion
\[
  x = \sum_{i=0}^{9} x_i 3^i,
  \qquad
  x_i\in\{0,1,2\}.
\]
The least-significant trit is $x_0$.

\section{Ternary rotation}

The Malbolge rotate instruction moves the least-significant trit to the most
significant position. Define
\[
  \rho(x)
    = \left\lfloor\frac{x}{3}\right\rfloor
      + (x\bmod 3)3^9.
\]
Equivalently, if $x$ has trits $(x_9,\ldots,x_1,x_0)$, then
\[
  \rho(x) = (x_0,x_9,\ldots,x_2,x_1)_3.
\]
The result remains in $\mathcal{W}$.

\section{Crazy operation}

Define the digitwise operation
\[
  \gamma : \{0,1,2\}\times\{0,1,2\}\to\{0,1,2\}
\]
by the table
\[
\begin{array}{c|ccc}
\gamma(d,a) & a=0 & a=1 & a=2 \\
\hline
 d=0 & 1 & 0 & 0 \\
 d=1 & 1 & 0 & 2 \\
 d=2 & 2 & 2 & 1
\end{array}
\]
where $d$ is the data-word trit and $a$ is the accumulator trit.

For words $d,a\in\mathcal{W}$, the classic crazy operation is
\[
  \operatorname{crazy}(d,a)
    = \sum_{i=0}^{9} \gamma(d_i,a_i)3^i.
\]
Because each output digit is ternary, the result is also in $\mathcal{W}$.

\section{Correspondence obligation}

An executable implementation may replace these definitions with lookup tables,
vector operations, GPU kernels, native instructions, or other optimized forms
only when the admitted domain is observably equivalent to these definitions.
Optimization cost does not weaken this semantic obligation.

\end{document}
